
\documentclass{article}
\usepackage{colortbl}
\usepackage{makecell}
\usepackage{multirow}
\usepackage{supertabular}

\begin{document}

\newcounter{utterance}

\centering \large Interaction Transcript for game `dond', experiment `coop\_en', episode 16 with Qwen3.5{-}9B.
\vspace{24pt}

{ \footnotesize  \setcounter{utterance}{1}
\setlength{\tabcolsep}{0pt}
\begin{supertabular}{c@{$\;$}|p{.15\linewidth}@{}p{.15\linewidth}p{.15\linewidth}p{.15\linewidth}p{.15\linewidth}p{.15\linewidth}}
    \# & $\;$A & \multicolumn{4}{c}{Game Master} & $\;\:$B\\
    \hline

    \theutterance \stepcounter{utterance}  
    & & \multicolumn{4}{p{0.6\linewidth}}{
        \cellcolor[rgb]{0.9,0.9,0.9}{
            \makecell[{{p{\linewidth}}}]{
                \texttt{\tiny{[P1$\langle$GM]}}
                \texttt{You are playing a cooperative negotiation game in which you have to agree on how to divide a set of items among you and another player.} \\
\\ 
\texttt{Rules:} \\
\texttt{(a) You and the other player are given a set of items. Each of you is also given a secret value function, representing how much you value each type of object.} \\
\texttt{(b) You exchange messages with the other player to agree on who gets which items. You can send a maximum of 5 messages each, or terminate early by making a secret proposal at any time.} \\
\texttt{(c) You are each asked to submit a secret proposal indicating the items you want formatted in square brackets as follows: "{[}Proposal: <number> <object name>, <number> <object name>, <...>{]}"} \\
\texttt{(d) If your proposals are complementary, i.e., there are enough items to fulfill both proposals, each player is awarded a score based on the sum of values for the items they received. Otherwise, both of you get zero points.} \\
\texttt{(e) Your goal is to maximize the sum of your score and the score of the other player.} \\
\\ 
\texttt{Let us start.} \\
\\ 
\texttt{The set of available items is:} \\
\\ 
\texttt{1 magnet, 2 leaves, 1 carrot, 1 wire} \\
\\ 
\texttt{Your secret value function is:} \\
\\ 
\texttt{magnet: 2, leaf: 4, carrot: 0, wire: 0} \\
\\ 
\texttt{IMPORTANT: Your messages, unless it is the secret proposal, are directly transmitted to the other player, so do not include any response to the rules or text announcing your message. To make a secret proposal, use the indicated format. Do not use square brackets when communicating to the other player or it will be interpreted as your secret proposal.} \\
            }
        }
    }
    & \\ \\

    \theutterance \stepcounter{utterance}  
    & \multicolumn{4}{p{0.6\linewidth}}{
        \cellcolor[rgb]{0.9,0.9,0.9}{
            \makecell[{{p{\linewidth}}}]{
                \texttt{\tiny{[P1$\rangle$GM]}}
                \texttt{I value the leaves most (4 points each), followed by the magnet (2 points). The carrot and wire are worthless to me (0 points). Since our goal is to maximize the *total* combined score, we should ensure the items go to the person who values them most.} \\
\\ 
\texttt{I will propose taking the two leaves, as they are my highest value items. I will leave the magnet, carrot, and wire open for the other player to claim, assuming they might value the magnet or the other items more than I do. This avoids conflict over the leaves and allows the other player to potentially take the magnet if they value it.} \\
\\ 
\texttt{Proposal:} \\
\texttt{I want the 2 leaves.} \\
            }
        }
    }
    & & \\ \\

    \theutterance \stepcounter{utterance}  
    & & & \multicolumn{2}{p{0.3\linewidth}}{
        \cellcolor[rgb]{0.9,0.9,0.9}{
            \makecell[{{p{\linewidth}}}]{
                \texttt{\tiny{[GM$|$GM]}}
                \texttt{continue} \\
            }
        }
    }
    & & \\ \\

    \theutterance \stepcounter{utterance}  
    & & & \multicolumn{2}{p{0.3\linewidth}}{
        \cellcolor[rgb]{0.9,0.9,0.9}{
            \makecell[{{p{\linewidth}}}]{
                \texttt{\tiny{[GM$|$GM]}}
                \texttt{I value the leaves most (4 points each), followed by the magnet (2 points). The carrot and wire are worthless to me (0 points). Since our goal is to maximize the *total* combined score, we should ensure the items go to the person who values them most.} \\
\\ 
\texttt{I will propose taking the two leaves, as they are my highest value items. I will leave the magnet, carrot, and wire open for the other player to claim, assuming they might value the magnet or the other items more than I do. This avoids conflict over the leaves and allows the other player to potentially take the magnet if they value it.} \\
\\ 
\texttt{Proposal:} \\
\texttt{I want the 2 leaves.} \\
            }
        }
    }
    & & \\ \\

    \theutterance \stepcounter{utterance}  
    & & \multicolumn{4}{p{0.6\linewidth}}{
        \cellcolor[rgb]{0.9,0.9,0.9}{
            \makecell[{{p{\linewidth}}}]{
                \texttt{\tiny{[GM$\rangle$P2]}}
                \texttt{You are playing a cooperative negotiation game in which you have to agree on how to divide a set of items among you and another player.} \\
\\ 
\texttt{Rules:} \\
\texttt{(a) You and the other player are given a set of items. Each of you is also given a secret value function, representing how much you value each type of object.} \\
\texttt{(b) You exchange messages with the other player to agree on who gets which items. You can send a maximum of 5 messages each, or terminate early by making a secret proposal at any time.} \\
\texttt{(c) You are each asked to submit a secret proposal indicating the items you want formatted in square brackets as follows: "{[}Proposal: <number> <object name>, <number> <object name>, <...>{]}"} \\
\texttt{(d) If your proposals are complementary, i.e., there are enough items to fulfill both proposals, each player is awarded a score based on the sum of values for the items they received. Otherwise, both of you get zero points.} \\
\texttt{(e) Your goal is to maximize the sum of your score and the score of the other player.} \\
\\ 
\texttt{Let us start.} \\
\\ 
\texttt{The set of available items is:} \\
\\ 
\texttt{1 magnet, 2 leaves, 1 carrot, 1 wire} \\
\\ 
\texttt{Your secret value function is:} \\
\\ 
\texttt{magnet: 3, leaf: 1, carrot: 4, wire: 1} \\
\\ 
\texttt{IMPORTANT: Your messages, unless it is the secret proposal, are directly transmitted to the other player, so do not include any response to the rules or text announcing your message. To make a secret proposal, use the indicated format. Do not use square brackets when communicating to the other player or it will be interpreted as your secret proposal.} \\
\\ 
\\ 
\texttt{I value the leaves most (4 points each), followed by the magnet (2 points). The carrot and wire are worthless to me (0 points). Since our goal is to maximize the *total* combined score, we should ensure the items go to the person who values them most.} \\
\\ 
\texttt{I will propose taking the two leaves, as they are my highest value items. I will leave the magnet, carrot, and wire open for the other player to claim, assuming they might value the magnet or the other items more than I do. This avoids conflict over the leaves and allows the other player to potentially take the magnet if they value it.} \\
\\ 
\texttt{Proposal:} \\
\texttt{I want the 2 leaves.} \\
            }
        }
    }
    & \\ \\

    \theutterance \stepcounter{utterance}  
    & & & \multicolumn{4}{p{0.6\linewidth}}{
        \cellcolor[rgb]{0.9,0.9,0.9}{
            \makecell[{{p{\linewidth}}}]{
                \texttt{\tiny{[GM$\langle$P2]}}
                \texttt{I understand the goal is to maximize the total combined score. Let's analyze the situation based on the information provided.} \\
\\ 
\texttt{**Item Analysis:**} \\
\texttt{*   **Leaves (2 available):** You value them at 4 each (Total 8). The other player values them at 4 each (Total 8).} \\
\texttt{*   **Magnet (1 available):** You value it at 3. The other player values it at 2.} \\
\texttt{*   **Carrot (1 available):** You value it at 4. The other player values it at 0.} \\
\texttt{*   **Wire (1 available):** You value it at 1. The other player values it at 0.} \\
\\ 
\texttt{**Strategy:**} \\
\texttt{Since the goal is the sum of our scores, items should go to whoever values them most.} \\
\texttt{1.  **Leaves:** Both of us value them equally (4). It doesn't matter who gets them for the *total* score, but we need to avoid conflict. If we both ask for leaves, we get 0. If one takes them, the total value is 8.} \\
\texttt{2.  **Carrot:** You value it at 4, the other player at 0. The carrot *must* go to you to contribute to the total score.} \\
\texttt{3.  **Magnet:** You value it at 3, the other player at 2. You value it more. It should go to you.} \\
\texttt{4.  **Wire:** You value it at 1, the other player at 0. You value it more. It should go to you.} \\
\\ 
\texttt{**Optimal Allocation for Max Total Score:**} \\
\texttt{*   **You:** 2 Leaves (8) + 1 Magnet (3) + 1 Carrot (4) + 1 Wire (1) = **16 points**.} \\
\texttt{*   **Other Player:** 0 points.} \\
\texttt{*   **Total Combined Score:** 16.} \\
\\ 
\texttt{If we split the leaves (e.g., you get 1, they get 1), the total leaf value is still 8. However, if the other player takes the magnet (valued at 2 by them, 3 by you), the total drops by 1. If they take the carrot (valued at 0 by them, 4 by you), the total drops by 4.} \\
\\ 
\texttt{Therefore, the most efficient proposal is for you to take everything, as you value every item more than or equal to the other player, and specifically strictly more for the magnet and carrot. There is no reason for the other player to take anything, as it would only reduce the total score if they take something you value more, or keep it the same if they take something you value equally (leaves).} \\
\\ 
\texttt{To ensure we don't conflict over the leaves (which are the only items we value equally), I will propose taking the leaves, the magnet, the} \\
            }
        }
    }
     \\ \\

    \theutterance \stepcounter{utterance}  
    & & & \multicolumn{2}{p{0.3\linewidth}}{
        \cellcolor[rgb]{0.9,0.9,0.9}{
            \makecell[{{p{\linewidth}}}]{
                \texttt{\tiny{[GM$|$GM]}}
                \texttt{continue} \\
            }
        }
    }
    & & \\ \\

    \theutterance \stepcounter{utterance}  
    & & & \multicolumn{2}{p{0.3\linewidth}}{
        \cellcolor[rgb]{0.9,0.9,0.9}{
            \makecell[{{p{\linewidth}}}]{
                \texttt{\tiny{[GM$|$GM]}}
                \texttt{I understand the goal is to maximize the total combined score. Let's analyze the situation based on the information provided.} \\
\\ 
\texttt{**Item Analysis:**} \\
\texttt{*   **Leaves (2 available):** You value them at 4 each (Total 8). The other player values them at 4 each (Total 8).} \\
\texttt{*   **Magnet (1 available):** You value it at 3. The other player values it at 2.} \\
\texttt{*   **Carrot (1 available):** You value it at 4. The other player values it at 0.} \\
\texttt{*   **Wire (1 available):** You value it at 1. The other player values it at 0.} \\
\\ 
\texttt{**Strategy:**} \\
\texttt{Since the goal is the sum of our scores, items should go to whoever values them most.} \\
\texttt{1.  **Leaves:** Both of us value them equally (4). It doesn't matter who gets them for the *total* score, but we need to avoid conflict. If we both ask for leaves, we get 0. If one takes them, the total value is 8.} \\
\texttt{2.  **Carrot:** You value it at 4, the other player at 0. The carrot *must* go to you to contribute to the total score.} \\
\texttt{3.  **Magnet:** You value it at 3, the other player at 2. You value it more. It should go to you.} \\
\texttt{4.  **Wire:** You value it at 1, the other player at 0. You value it more. It should go to you.} \\
\\ 
\texttt{**Optimal Allocation for Max Total Score:**} \\
\texttt{*   **You:** 2 Leaves (8) + 1 Magnet (3) + 1 Carrot (4) + 1 Wire (1) = **16 points**.} \\
\texttt{*   **Other Player:** 0 points.} \\
\texttt{*   **Total Combined Score:** 16.} \\
\\ 
\texttt{If we split the leaves (e.g., you get 1, they get 1), the total leaf value is still 8. However, if the other player takes the magnet (valued at 2 by them, 3 by you), the total drops by 1. If they take the carrot (valued at 0 by them, 4 by you), the total drops by 4.} \\
\\ 
\texttt{Therefore, the most efficient proposal is for you to take everything, as you value every item more than or equal to the other player, and specifically strictly more for the magnet and carrot. There is no reason for the other player to take anything, as it would only reduce the total score if they take something you value more, or keep it the same if they take something you value equally (leaves).} \\
\\ 
\texttt{To ensure we don't conflict over the leaves (which are the only items we value equally), I will propose taking the leaves, the magnet, the} \\
            }
        }
    }
    & & \\ \\

    \theutterance \stepcounter{utterance}  
    & & \multicolumn{4}{p{0.6\linewidth}}{
        \cellcolor[rgb]{0.9,0.9,0.9}{
            \makecell[{{p{\linewidth}}}]{
                \texttt{\tiny{[P1$\langle$GM]}}
                \texttt{I understand the goal is to maximize the total combined score. Let's analyze the situation based on the information provided.} \\
\\ 
\texttt{**Item Analysis:**} \\
\texttt{*   **Leaves (2 available):** You value them at 4 each (Total 8). The other player values them at 4 each (Total 8).} \\
\texttt{*   **Magnet (1 available):** You value it at 3. The other player values it at 2.} \\
\texttt{*   **Carrot (1 available):** You value it at 4. The other player values it at 0.} \\
\texttt{*   **Wire (1 available):** You value it at 1. The other player values it at 0.} \\
\\ 
\texttt{**Strategy:**} \\
\texttt{Since the goal is the sum of our scores, items should go to whoever values them most.} \\
\texttt{1.  **Leaves:** Both of us value them equally (4). It doesn't matter who gets them for the *total* score, but we need to avoid conflict. If we both ask for leaves, we get 0. If one takes them, the total value is 8.} \\
\texttt{2.  **Carrot:** You value it at 4, the other player at 0. The carrot *must* go to you to contribute to the total score.} \\
\texttt{3.  **Magnet:** You value it at 3, the other player at 2. You value it more. It should go to you.} \\
\texttt{4.  **Wire:** You value it at 1, the other player at 0. You value it more. It should go to you.} \\
\\ 
\texttt{**Optimal Allocation for Max Total Score:**} \\
\texttt{*   **You:** 2 Leaves (8) + 1 Magnet (3) + 1 Carrot (4) + 1 Wire (1) = **16 points**.} \\
\texttt{*   **Other Player:** 0 points.} \\
\texttt{*   **Total Combined Score:** 16.} \\
\\ 
\texttt{If we split the leaves (e.g., you get 1, they get 1), the total leaf value is still 8. However, if the other player takes the magnet (valued at 2 by them, 3 by you), the total drops by 1. If they take the carrot (valued at 0 by them, 4 by you), the total drops by 4.} \\
\\ 
\texttt{Therefore, the most efficient proposal is for you to take everything, as you value every item more than or equal to the other player, and specifically strictly more for the magnet and carrot. There is no reason for the other player to take anything, as it would only reduce the total score if they take something you value more, or keep it the same if they take something you value equally (leaves).} \\
\\ 
\texttt{To ensure we don't conflict over the leaves (which are the only items we value equally), I will propose taking the leaves, the magnet, the} \\
            }
        }
    }
    & \\ \\

    \theutterance \stepcounter{utterance}  
    & \multicolumn{4}{p{0.6\linewidth}}{
        \cellcolor[rgb]{0.9,0.9,0.9}{
            \makecell[{{p{\linewidth}}}]{
                \texttt{\tiny{[P1$\rangle$GM]}}
                \texttt{Proposal: 2 leaf, 1 magnet, 1 carrot, 1 wire} \\
            }
        }
    }
    & & \\ \\

    \theutterance \stepcounter{utterance}  
    & & & \multicolumn{2}{p{0.3\linewidth}}{
        \cellcolor[rgb]{0.9,0.9,0.9}{
            \makecell[{{p{\linewidth}}}]{
                \texttt{\tiny{[GM$|$GM]}}
                \texttt{continue} \\
            }
        }
    }
    & & \\ \\

    \theutterance \stepcounter{utterance}  
    & & & \multicolumn{2}{p{0.3\linewidth}}{
        \cellcolor[rgb]{0.9,0.9,0.9}{
            \makecell[{{p{\linewidth}}}]{
                \texttt{\tiny{[GM$|$GM]}}
                \texttt{Proposal: 2 leaf, 1 magnet, 1 carrot, 1 wire} \\
            }
        }
    }
    & & \\ \\

    \theutterance \stepcounter{utterance}  
    & & \multicolumn{4}{p{0.6\linewidth}}{
        \cellcolor[rgb]{0.9,0.9,0.9}{
            \makecell[{{p{\linewidth}}}]{
                \texttt{\tiny{[GM$\rangle$P2]}}
                \texttt{Proposal: 2 leaf, 1 magnet, 1 carrot, 1 wire} \\
            }
        }
    }
    & \\ \\

    \theutterance \stepcounter{utterance}  
    & & & \multicolumn{4}{p{0.6\linewidth}}{
        \cellcolor[rgb]{0.9,0.9,0.9}{
            \makecell[{{p{\linewidth}}}]{
                \texttt{\tiny{[GM$\langle$P2]}}
                \texttt{I see you have made a secret proposal. Since I am the other player in this simulation, I need to formulate my own secret proposal to complete the negotiation.} \\
\\ 
\texttt{**My Analysis:**} \\
\texttt{*   **My Values:** Leaf (4), Magnet (2), Carrot (0), Wire (0).} \\
\texttt{*   **Your Proposal:** 2 leaves, 1 magnet, 1 carrot, 1 wire.} \\
\texttt{*   **Conflict Check:** You are asking for *all* items. If I also ask for any item, our proposals will overlap, resulting in a "complementary failure" (both get 0 points).} \\
\texttt{*   **Optimal Strategy:** To maximize the total score, the items must go to the person who values them most.} \\
\texttt{    *   **Leaves:** You value 4, I value 4. (Tie).} \\
\texttt{    *   **Magnet:** You value 3, I value 2. (You win).} \\
\texttt{    *   **Carrot:** You value 4, I value 0. (You win).} \\
\texttt{    *   **Wire:** You value 1, I value 0. (You win).} \\
\\ 
\texttt{Since you are claiming every single item, and you value every item at least as much as I do (with strict superiority on magnet, carrot, and wire), the only way to achieve a non{-}zero score is for me to claim **nothing**. If I claim anything, we clash and both get 0. If I claim nothing, you get everything, and the total score is maximized based on your valuations.} \\
\\ 
\texttt{**My Secret Proposal:**} \\
\texttt{{[}Proposal: 0 magnet, 0 leaf, 0 carrot, 0 wire{]}} \\
\\ 
\texttt{*(Note: In the context of the game rules provided, submitting a proposal with 0 quantities effectively means "I want nothing," which is complementary to your proposal of taking everything.)*} \\
            }
        }
    }
     \\ \\

    \theutterance \stepcounter{utterance}  
    & & & \multicolumn{2}{p{0.3\linewidth}}{
        \cellcolor[rgb]{0.9,0.9,0.9}{
            \makecell[{{p{\linewidth}}}]{
                \texttt{\tiny{[GM$|$GM]}}
                \texttt{proposal} \\
            }
        }
    }
    & & \\ \\

    \theutterance \stepcounter{utterance}  
    & & & \multicolumn{2}{p{0.3\linewidth}}{
        \cellcolor[rgb]{0.9,0.9,0.9}{
            \makecell[{{p{\linewidth}}}]{
                \texttt{\tiny{[GM$|$GM]}}
                [0, 0, 0, 0]
            }
        }
    }
    & & \\ \\

    \theutterance \stepcounter{utterance}  
    & & \multicolumn{4}{p{0.6\linewidth}}{
        \cellcolor[rgb]{0.9,0.9,0.9}{
            \makecell[{{p{\linewidth}}}]{
                \texttt{\tiny{[P1$\langle$GM]}}
                \texttt{The other player has submitted a secret proposal. It is now your turn to submit a proposal of your own.} \\
            }
        }
    }
    & \\ \\

    \theutterance \stepcounter{utterance}  
    & \multicolumn{4}{p{0.6\linewidth}}{
        \cellcolor[rgb]{0.9,0.9,0.9}{
            \makecell[{{p{\linewidth}}}]{
                \texttt{\tiny{[P1$\rangle$GM]}}
                \texttt{Proposal: 2 leaf, 1 magnet, 1 carrot, 1 wire} \\
            }
        }
    }
    & & \\ \\

    \theutterance \stepcounter{utterance}  
    & & & \multicolumn{2}{p{0.3\linewidth}}{
        \cellcolor[rgb]{0.9,0.9,0.9}{
            \makecell[{{p{\linewidth}}}]{
                \texttt{\tiny{[GM$|$GM]}}
                \texttt{continue} \\
            }
        }
    }
    & & \\ \\

    \theutterance \stepcounter{utterance}  
    & & & \multicolumn{2}{p{0.3\linewidth}}{
        \cellcolor[rgb]{0.9,0.9,0.9}{
            \makecell[{{p{\linewidth}}}]{
                \texttt{\tiny{[GM$|$GM]}}
                \texttt{player was instructed to make a proposal} \\
            }
        }
    }
    & & \\ \\

\end{supertabular}
}

\end{document}
